\documentclass[aspectratio=169,11pt,spanish]{beamer}
\input{../../../_preambulo_beamer.tex}
\input{../../../_comandos_beamer.tex}
\title[L05 · Fundamentos]{Teoría de conjuntos y espacios muestrales}
\begin{document}
\shorthandoff{"}
\begin{frame}{L05 · Un experimento y sus resultados}
Un \textbf{experimento aleatorio} tiene un resultado que no conocemos antes de realizarlo.
Consideremos un lanzamiento de un dado de seis caras.
\[
S=\{1,2,3,4,5,6\},\qquad |S|=6.
\]
El \textbf{espacio muestral} $S$ contiene todos los resultados posibles. En un lanzamiento se observa un elemento de $S$.

\medskip
La \textbf{cardinalidad} $|S|$ es el número de elementos distintos de $S$.
\end{frame}
\begin{frame}{Un evento es un subconjunto}
Un \textbf{evento} reúne los resultados que cumplen una condición. Para el dado, definimos:
\[
A=\{2,4,6\}\quad\text{(sale un número par)},\qquad
B=\{4,5,6\}\quad\text{(sale al menos cuatro)}.
\]
Ambos son subconjuntos de $S$ y $|A|=|B|=3$.

\medskip
Si sale 4, ocurren $A$ y $B$: un resultado puede pertenecer a más de un evento.
\end{frame}
\begin{frame}{Intersección: cumplir ambas condiciones}
La intersección conserva sólo los elementos comunes:
\[
A\cap B=\{x\in S:x\in A\text{ y }x\in B\}.
\]
En nuestro ejemplo,
\[
A\cap B=\{4,6\},\qquad |A\cap B|=2.
\]
El 4 y el 6 son pares y, a la vez, son al menos cuatro.
\end{frame}
\begin{frame}{Unión: cumplir al menos una condición}
La unión reúne los elementos de $A$, de $B$ o de ambos:
\[
A\cup B=\{x\in S:x\in A\text{ o }x\in B\}.
\]
La \textbf{o} es inclusiva. Por tanto,
\[
A\cup B=\{2,4,5,6\},\qquad |A\cup B|=4.
\]
Cada resultado se escribe una sola vez en un conjunto.
\end{frame}
\begin{frame}{Contar una unión sin duplicar elementos}
Al sumar $|A|+|B|$, los elementos de $A\cap B$ se cuentan dos veces. Restamos una copia de cada uno:
\[
|A\cup B|=|A|+|B|-|A\cap B|.
\]
Para los eventos definidos antes:
\[
4=3+3-2.
\]
Esta identidad es válida para cualesquiera dos conjuntos finitos.
\end{frame}
\begin{frame}{Complemento: respecto a un universo}
El complemento de $A$ contiene los resultados de $S$ que no están en $A$:
\[
A^c=S\setminus A=\{1,3,5\},\qquad |A^c|=3.
\]
$A$ y $A^c$ separan todo el espacio muestral:
\[
A\cap A^c=\varnothing,\quad A\cup A^c=S,\quad
|A^c|=|S|-|A|=6-3=3.
\]
Sin un universo definido, el complemento queda ambiguo.
\end{frame}
\begin{frame}{Disyunción y exhaustividad}
Dos eventos $X,Y\subseteq S$ son \textbf{disjuntos} si $X\cap Y=\varnothing$.
Son \textbf{exhaustivos} si $X\cup Y=S$.

\medskip
Sean $C=\{1,3,5\}$ (impar) y $D=\{2,4,6\}$ (par). Entonces
\[
C\cap D=\varnothing,\qquad C\cup D=S.
\]
$C$ y $D$ tienen ambas propiedades. Disyunción pregunta por el \textbf{traslape}; exhaustividad pregunta por la \textbf{cobertura}.
\end{frame}
\begin{frame}{Ideas para comprobar tus resultados}
Para eventos $A,B\subseteq S$:
\[
A\cap B\subseteq A,\quad A\cap B\subseteq B,\qquad
A\subseteq A\cup B,\quad B\subseteq A\cup B.
\]
Por ello $|A\cap B|$ no puede superar $|A|$ o $|B|$, y $|A\cup B|$ no puede ser menor que ellos.

\medskip
\scriptsize Referencia: OpenStax, \emph{Introductory Business Statistics 2e}, secciones 3.1 y 3.2. El ejemplo del dado es autónomo.
\end{frame}
\end{document}
